% Prose rendering of PrimeTensor/Depth.lean.
% This file is a fragment: it carries no \documentclass and no document
% environment, so it can be \input into any document that loads
% proof/preamble.tex.  See proof/README.md.

\section{Positive depth and nonempty finite axes}
\label{sec:depth}

\leanfile{PrimeTensor/Depth.lean}

\subsection*{What this file establishes}

The file fixes the indexing vocabulary that the rest of the development is
built on: a type of \emph{depths} (dimensions and ranks), a family of
coordinate \emph{axes} indexed by a depth, a fold over an axis, and the type
of index tuples of a given rank.

Its single organising decision is that \emph{nothing starts at zero}.
$\Depth$ has no zeroth constructor; consequently every axis carries at least
one coordinate, every tensor has at least one dimension, and every index
tuple has at least one component.  This is what lets $\fold$ be defined
without an identity element for its binary operation, which in turn is what
makes the later multiplicative layer --- where the carrier is the strictly
positive reals and there is no additive unit to fall back on --- expressible
at all.

\subsection{Depth}

\begin{definition}[Depth, \lean{PrimeTensor.Depth}]\label{def:depth}
$\Depth$ is the inductive type generated by the two constructors
\[
  \one : \Depth,
  \qquad
  \suc : \Depth \to \Depth .
\]
\end{definition}

\begin{lstlisting}
inductive Depth where
  | one : Depth
  | succ : Depth → Depth
  deriving DecidableEq, Repr
\end{lstlisting}

This is the Peano presentation of the natural numbers with the base case
moved from $0$ to $1$.  Writing
\[
  \size{\one} \defeq 1,
  \qquad
  \size{\suc d} \defeq \size{d} + 1,
\]
the map $d \mapsto \size{d}$ is a bijection $\Depth \to \mathbb{Z}_{\geq 1}$;
we use it silently below whenever a depth is read as a number.  Because
$\Depth$ is a first-order inductive type with no parameters, Lean derives
decidable equality and a printing instance for it mechanically, and structural
induction on $\Depth$ is available: to prove $P(d)$ for all $d$ it suffices to
prove $P(\one)$ and $P(d) \Rightarrow P(\suc d)$.

\begin{definition}[Small depths, \lean{Depth.two}, \lean{Depth.three}]
\label{def:small-depths}
\[
  \mathsf{two} \defeq \suc\one,
  \qquad
  \mathsf{three} \defeq \suc\mathsf{two} .
\]
\end{definition}

These are abbreviations, not new data: $\size{\mathsf{two}} = 2$ and
$\size{\mathsf{three}} = 3$, and each unfolds definitionally to a stack of
constructors.  \textsf{three} is the depth at which the Navier--Stokes layer
of the project operates.

\begin{designnote}\label{note:no-zero}
The absence of a zero constructor is a positive design choice and not an
oversight.  Had $\Depth$ been $\mathbb{N}$, then $\Axis\,0$
(Definition~\ref{def:axis}) would be empty, $\fold$
(Definition~\ref{def:fold}) would need an identity element to return in that
case, and every downstream lemma about a fold would carry a side condition
excluding the degenerate dimension.  Ruling the degenerate case out at the
level of the index type discharges all of those obligations at once, by making
the bad case unrepresentable rather than merely excluded.
\end{designnote}

\subsection{Axes}

\begin{definition}[Axis, \lean{PrimeTensor.Axis}]\label{def:axis}
$\Axis$ is the inductive family $\Depth \to \mathrm{Type}$ generated by
\[
  \first : \Axis\,d \quad (\text{for every } d : \Depth),
  \qquad
  \next : \Axis\,d \to \Axis\,(\suc d).
\]
\end{definition}

\begin{lstlisting}
inductive Axis : Depth → Type where
  | first {d : Depth} : Axis d
  | next  {d : Depth} : Axis d → Axis (.succ d)
\end{lstlisting}

Note the asymmetry between the two constructors.  $\first$ is available at
every depth, so it is the coordinate that always exists; $\next$ is available
only at a depth that is literally a successor, so it can only push an existing
coordinate of $\Axis\,d$ into the strictly larger $\Axis\,(\suc d)$.

\begin{proposition}[Axes are inhabited]\label{prop:axis-inhabited}
For every $d : \Depth$ the type $\Axis\,d$ is nonempty.
\end{proposition}

\begin{proof}
The term $\first : \Axis\,d$ is well formed for every $d$, since the depth
argument of $\first$ is unconstrained.
\end{proof}

\begin{proposition}[Cardinality of an axis]\label{prop:axis-card}
For every $d : \Depth$ the type $\Axis\,d$ is finite, with exactly
$\size{d}$ elements.
\end{proposition}

\begin{proof}
By structural induction on $d$ (Definition~\ref{def:depth}).

\emph{Base case $d = \one$.}  Every element of an inductive family is a
constructor application.  An element of $\Axis\,\one$ built by $\next$ would
have type $\Axis\,(\suc d')$ for some $d'$, so its depth index would have to
satisfy $\one = \suc d'$; distinct constructors of $\Depth$ produce distinct
values, so no such $d'$ exists.  Hence $\first$ is the only element, and
$\#\Axis\,\one = 1 = \size{\one}$.

\emph{Inductive step $d = \suc d'$.}  The two constructors are jointly
surjective onto $\Axis\,(\suc d')$ and have disjoint images, and $\next$ is
injective, both again because constructors of an inductive type are disjoint
and injective.  Therefore
\[
  \Axis\,(\suc d') \;\cong\; \{\first\} \sqcup \Axis\,d',
\]
so by the induction hypothesis
$\#\Axis\,(\suc d') = 1 + \size{d'} = \size{\suc d'}$.
\end{proof}

\begin{notation}\label{not:axis-enum}
Proposition~\ref{prop:axis-card} lets us name the elements of $\Axis\,d$ in
order.  For $n = \size{d}$ put
\[
  i_1 \defeq \first,
  \qquad
  i_{k+1} \defeq \next\,i_k \quad (1 \leq k < n),
\]
so that $\Axis\,d = \{i_1, \dots, i_n\}$ with all $i_k$ distinct.  Equivalently
$i_k = \next^{\,k-1}(\first)$, and this term is well typed in $\Axis\,d$
exactly when $k \leq n$: each application of $\next$ consumes one $\sucbare$ from
the depth index, so $d$ must be at least a $(k-1)$-fold successor.  The
enumeration is thus cut off by the depth itself, with no side condition to
carry.
\end{notation}

\subsection{Folding over an axis}

\begin{definition}[Fold, \lean{PrimeTensor.Axis.fold}]\label{def:fold}
Let $A$ be a type and $\ast : A \to A \to A$ a binary operation.  For
$d : \Depth$ and $f : \Axis\,d \to A$ define $\fold_{\ast}(d, f) : A$ by
structural recursion on $d$:
\begin{align}
  \fold_{\ast}(\one, f)
    &= f(\first), \label{eq:fold-one}\\
  \fold_{\ast}(\suc d, f)
    &= f(\first) \ast \fold_{\ast}\bigl(d,\; f \circ \next\bigr).
      \label{eq:fold-succ}
\end{align}
\end{definition}

\begin{lstlisting}
def fold {A : Type} (op : A → A → A) :
    (d : Depth) → (Axis d → A) → A
  | .one,    f => f .first
  | .succ d, f => op (f .first) (fold op d (fun i => f (.next i)))
\end{lstlisting}

The recursion is structural in $d$ and therefore terminates: the recursive
call is made at the immediate predecessor of $\suc d$.  The function argument
is re-indexed rather than shrunk --- $f \circ \next$ has domain $\Axis\,d$
because $\next : \Axis\,d \to \Axis\,(\suc d)$ --- which is exactly why the
recursion can be driven by the depth alone.

\begin{designnote}\label{note:no-identity}
$\fold$ takes no seed value and $\ast$ is required to be neither associative,
nor commutative, nor unital.  A fold over a possibly-empty index type must
return something on the empty index type, and the only uniform choice is an
identity element for $\ast$; requiring one would force every carrier used with
$\fold$ to be a monoid.  Proposition~\ref{prop:axis-inhabited} removes the
empty case, so $\fold$ is total on a bare magma.  This is what makes $\fold$
usable over the strictly positive reals under multiplication, and over
$\min$ and $\max$, none of which is a monoid in the form the later files need.
\end{designnote}

\begin{lemma}[Base equation, \lean{PrimeTensor.Axis.fold\_one}]
\label{lem:fold-one}
For all $\ast$ and all $f : \Axis\,\one \to A$,
\[
  \fold_{\ast}(\one, f) = f(\first).
\]
\end{lemma}

\begin{proof}
Both sides are the same term after unfolding $\fold$ once: \eqref{eq:fold-one}
is the defining equation of $\fold$ at the constructor $\one$, and the match
on $d$ reduces because $\one$ is a constructor application in head normal
form.  The two sides are therefore definitionally equal, and the proof is
$\mathsf{rfl}$.
\end{proof}

\begin{lemma}[Successor equation, \lean{PrimeTensor.Axis.fold\_succ}]
\label{lem:fold-succ}
For all $\ast$, all $d : \Depth$, and all $f : \Axis\,(\suc d) \to A$,
\[
  \fold_{\ast}(\suc d, f)
    = f(\first) \ast \fold_{\ast}\bigl(d,\; f \circ \next\bigr).
\]
\end{lemma}

\begin{proof}
As in Lemma~\ref{lem:fold-one}: this is \eqref{eq:fold-succ}, the defining
equation at the constructor $\sucbare$, and the match reduces without any
hypothesis on $d$.  Definitional equality again, proved by $\mathsf{rfl}$.
\end{proof}

\begin{remark}\label{rem:simp}
Both equations hold by reflexivity, so as \emph{propositions} they are
content-free.  They are stated, and tagged \lean{@[simp]}, for a mechanical
reason: Lean's simplifier rewrites with declared lemmas, not by unfolding
recursive definitions on its own.  Registering the two equations makes
$\mathsf{simp}$ evaluate any fold whose depth is a concrete numeral, which is
what downstream files rely on.  Together they also constitute the complete
case analysis on $d$, so no third equation is possible.
\end{remark}

\begin{proposition}[What a fold computes]\label{prop:fold-expansion}
Let $n = \size{d}$ and let $i_1, \dots, i_n$ enumerate $\Axis\,d$ as in
Notation~\ref{not:axis-enum}.  Then
\[
  \fold_{\ast}(d, f)
    = f(i_1) \ast \bigl( f(i_2) \ast \bigl( \cdots \ast f(i_n) \bigr) \bigr),
\]
the right-nested $\ast$-expression over the coordinates in order, with no
identity element at the innermost position.
\end{proposition}

\begin{proof}
By induction on $d$.  For $d = \one$ we have $n = 1$ and
Lemma~\ref{lem:fold-one} gives $f(i_1)$.  For $d = \suc d'$,
Lemma~\ref{lem:fold-succ} gives
$f(\first) \ast \fold_{\ast}(d', f \circ \next)$; the enumeration of
$\Axis\,(\suc d')$ is $\first$ followed by the $\next$-image of the
enumeration of $\Axis\,d'$, so the induction hypothesis applied to
$f \circ \next$ rewrites the second factor as the right-nested expression
over $f(i_2), \dots, f(i_n)$.
\end{proof}

\begin{remark}
Proposition~\ref{prop:fold-expansion} is the informal reading of
Definition~\ref{def:fold}; it is stated here for the reader and is not among
the declarations of \texttt{PrimeTensor/Depth.lean}.  The bracketing it
exhibits is fixed, not incidental: without associativity of $\ast$ the value
of a fold depends on it.
\end{remark}

\subsection{Index tuples}

\begin{definition}[Index tuple, \lean{PrimeTensor.IndexTuple}]
\label{def:indextuple}
For a dimension $\mathit{dim} : \Depth$ define $\Idx(\mathit{dim}, -)$, a
map $\Depth \to \mathrm{Type}$, by structural recursion on the rank:
\begin{align*}
  \Idx(\mathit{dim}, \one)     &= \Axis\,\mathit{dim}, \\
  \Idx(\mathit{dim}, \suc r)   &= \Axis\,\mathit{dim} \times
                                   \Idx(\mathit{dim}, r).
\end{align*}
\end{definition}

\begin{lstlisting}
def IndexTuple (dim : Depth) : Depth → Type
  | .one   => Axis dim
  | .succ r => Axis dim × IndexTuple dim r
\end{lstlisting}

Unlike $\Axis$, this is not an inductive family but a function into
$\mathrm{Type}$ defined by recursion, so its two equations hold
definitionally and a tuple is literally a nested pair --- no constructor or
projection layer sits between $\Idx(\mathit{dim}, \mathsf{three})$ and
$\Axis\,\mathit{dim} \times (\Axis\,\mathit{dim} \times \Axis\,\mathit{dim})$.
Both arguments are depths, but they play different roles: $\mathit{dim}$ is
the size of each coordinate axis and is held fixed, while the second argument
is the rank and drives the recursion.  A rank-one index tuple is a single axis
index rather than a one-element product, so there is no empty tuple and no
unit type in the encoding.

\begin{proposition}[Cardinality of the index set]\label{prop:idx-card}
$\Idx(\mathit{dim}, r)$ is finite with
$\size{\mathit{dim}}^{\,\size{r}}$ elements; in particular it is nonempty.
\end{proposition}

\begin{proof}
By induction on $r$.  For $r = \one$, Proposition~\ref{prop:axis-card} gives
$\#\Axis\,\mathit{dim} = \size{\mathit{dim}} = \size{\mathit{dim}}^{1}$.  For
$r = \suc r'$, the cardinality of a product is the product of the
cardinalities, so
\[
  \#\Idx(\mathit{dim}, \suc r')
    = \size{\mathit{dim}} \cdot \size{\mathit{dim}}^{\,\size{r'}}
    = \size{\mathit{dim}}^{\,\size{r'} + 1}
    = \size{\mathit{dim}}^{\,\size{\suc r'}}.
\]
Nonemptiness follows since $\size{\mathit{dim}} \geq 1$, or directly from
Proposition~\ref{prop:axis-inhabited} by pairing.
\end{proof}

\begin{corollary}[The intended case]\label{cor:rank-three}
$\Idx(\mathsf{three}, \mathsf{three})$ has $3^3 = 27$ elements: the index set
of a rank-three tensor in three dimensions.
\end{corollary}

\subsection*{Summary of the correspondence}

\begin{center}
\small
\begin{tabular}{@{}lll@{}}
\hline
Lean declaration & Kind & Here \\
\hline
\lean{PrimeTensor.Depth}          & inductive type   & Def.~\ref{def:depth} \\
\lean{PrimeTensor.Depth.two}      & definition       & Def.~\ref{def:small-depths} \\
\lean{PrimeTensor.Depth.three}    & definition       & Def.~\ref{def:small-depths} \\
\lean{PrimeTensor.Axis}           & inductive family & Def.~\ref{def:axis} \\
\lean{PrimeTensor.Axis.fold}      & definition       & Def.~\ref{def:fold} \\
\lean{PrimeTensor.Axis.fold\_one} & theorem          & Lem.~\ref{lem:fold-one} \\
\lean{PrimeTensor.Axis.fold\_succ}& theorem          & Lem.~\ref{lem:fold-succ} \\
\lean{PrimeTensor.IndexTuple}     & definition       & Def.~\ref{def:indextuple} \\
\hline
\end{tabular}
\end{center}

\noindent
Every declaration of \texttt{PrimeTensor/Depth.lean} appears above.  The file
contains no \lean{sorry}, no axiom, and no unproved named proposition; the two
theorems it does state are proved by $\mathsf{rfl}$.  The mathematical content
of the file is therefore entirely in its definitions, and the results proved
about them here --- Propositions~\ref{prop:axis-inhabited},
\ref{prop:axis-card}, \ref{prop:fold-expansion} and~\ref{prop:idx-card} ---
are consequences of those definitions stated for the reader rather than
formalised obligations.
